\section{Data and Evaluation Protocol}\label{sec:data}

\subsection{EgoPoint}
We train and test on the image set of EgoPoint-Ground~\cite{egopoint}, which we call EgoPoint. Its paper reports 15,338 images: 39\% (6,035) real captures with RayNeo~V3 smart glasses (12\,MP Sony IMX681 sensor), 39\% (5,976) synthetically composed and 22\% (3,327) LLM-generated; its appendix calls the last two groups ``synthetic generation'' and ``image editing''. The edited images were made from third-person photographs of COCO~\cite{lin2014coco}, selected through the RefCOCO referring-expression dataset~\cite{yu2016modeling}, into which a multimodal model added a pointing hand~\cite{egopoint}. The release we use has 15,455 images split 10,837 / 1,546 / 3,072. VistaRef~\cite{li2026vistaref} and PointVG-R~\cite{li2026pointvgr} report the same counts, so our test images are theirs, minus the exclusions below. The release does not mark which images are real captures and which are generated, so, like VistaRef and PointVG-R, we report on the whole test split. It does, however, name the source photograph of each edited image. In our test split, 842 of the 2,604 images with a target (32\%) are edited COCO train2014 photographs, and at least 400 of the 425 no-hand images are unedited COCO train2014 photographs, so at least 1,242 test images (41\%) come from COCO. This matters for baselines pretrained on COCO (Section~\ref{sec:metrics}). Each positive record shows one wearer's hand pointing at one target; negatives show no hand.

\paragraph{Conversion.} Negatives become 000, and their object boxes are ignored; positives with a hand box, an annotated fingertip and a target box become 111. In three test records, the hand label was attached to an object and the hand was labelled as an object; we detected these by the fingertip position and swapped the labels back. We dropped 232 positives: 229 without an annotated fingertip (164 train, 22 validation, 43 test), all flagged as unclear in the source annotations, and 3 without a hand box. Spot checks showed open palms, held objects, or a hand box on the hand that was not pointing. If all 43 dropped test positives were counted as failures, the target P@0.5 reported in Section~\ref{sec:results} would fall by less than one point, to about 58\%. Table~\ref{tab:data} summarises the result. Images are stored at $224\times224$; the model squashes every frame to its square $160\times160$ input, regardless of the original aspect ratio.

\begin{table}[t]
  \caption{EgoPoint after conversion. Gap categories (Section~\ref{sec:problem}): \emph{overlap}, the fingertip lies inside the target box; \emph{near}, fingertip--box distance ${\leq}\,0.08$ of the image side; \emph{large gap}, ${>}\,0.08$.}
  \label{tab:data}
  \footnotesize
  \begin{tabular}{lrrrrrr}
    \toprule
    & & & & \multicolumn{3}{c}{Pointing at target, by gap} \\
    \cmidrule(l){5-7}
    Split & Images & No hand (000) & Pointing at target (111) & Overlap & Near & Large gap \\
    \midrule
    Train & 10,669 & 1,535 & 9,134 & 6,511 & 1,995 & 628 \\
    Validation & 1,524 & 211 & 1,313 & 922 & 297 & 94 \\
    Test & 3,029 & 425 & 2,604 & 1,796 & 599 & 209 \\
    \midrule
    Total & 15,222 & 2,171 & 13,051 & 9,229 & 2,891 & 931 \\
    \bottomrule
  \end{tabular}
\end{table}

\paragraph{Consequences for evaluation.} Our test split has 3,029 images, 43 positives fewer than the official one, and published numbers on the official split are not strictly comparable with ours (Section~\ref{sec:main}). EgoPoint contains only scenes 000 and 111, so the pointing and target gates never see a negative example, and we report neither their accuracy nor four-class scene accuracy. About 70\% of targets overlap the fingertip and only 7\% have a large gap. Pooled results are dominated by overlap targets, so we report every target metric by gap as well. The median distance from the fingertip to the target centre is 0.12, 0.17 and 0.32 of the image side for overlap, near and large-gap targets (test split).

\subsection{Metrics}\label{sec:metrics}
\begin{itemize}
\item \emph{Target and hand localisation.} We use the two standard grounding metrics~\cite{egopoint,li2026vistaref,li2026pointvgr}: \emph{mIoU}, the IoU between predicted and labelled box averaged over images, with a 95\% bootstrap confidence interval (1,000 resamples~\cite{efron1994introduction}), and \emph{P@$\tau$}, the share of images with IoU${\geq}\tau$, for $\tau$ = 0.3, 0.5 and 0.7. Both are given in \%, and target metrics are given overall and by gap.
\item \emph{Fingertip and direction.} Fingertip error is in pixels at the $160\times160$ input (1\,px = 0.63\% of the image side). Direction error is the angle between $\hat{\dir}$ and $\dir^\star$.
\item \emph{Target choice.} \emph{Centre in target} is the share of target images whose predicted target centre falls inside the labelled box. It is a diagnostic of whether the model picked the right object, independent of how well it drew the box.
\item \emph{Hand presence.} For the hand gate, we report accuracy, missed hands, false alarms and AUROC, with Wilson intervals~\cite{wilson1927probable} for proportions.
\end{itemize}
Localisation is scored on every image whose label contains the part, whatever scene the model predicted, so hard images cannot be skipped.

\paragraph{End-to-end score.} The target metrics ignore the gates. To count missed hands and false alarms as well, we also report an \emph{end-to-end score} over all 3,029 test images. A target image counts as correct if the hand gate fires, the model predicts ``pointing at target'' and the target box has IoU${\geq}0.5$; a no-hand image counts as correct if the model says ``no hand''. For this score only, each model's hand-gate threshold is tuned on the validation split by the same rule.

\paragraph{Ownership test.} To test robustness to other people's hands, we paste one or two hands into every test image. The hands come from other test images, 30\% are re-pasted copies of the wearer's own hand, and their positions are fixed by a seed; labels stay unchanged. We report \emph{ownership accuracy} (the share of hand images in which the predicted hand centre is nearer the wearer's hand than any pasted hand), how often the predicted centre falls inside a pasted hand, and the share of no-hand images that stay \emph{no hand}. These distractors follow the training distribution, so this test measures robustness to that synthetic distribution and may overstate robustness in real scenes.

\paragraph{Compact detectors at matched compute.} The published baselines are large models that we did not re-run. To compare SpatialNet with models of its own size under exactly our protocol, we trained two compact detectors on the same EgoPoint training split: YOLO26n~\cite{ultralytics2026yolo26} at a $128\times128$ input, which matches SpatialNet's compute (0.225 against 0.231 GFLOPs), from three starting points (COCO-pretrained~\cite{lin2014coco}; ImageNet-pretrained, the same pretraining as our backbone~\cite{deng2009imagenet}; and from scratch), plus COCO-pretrained runs at 160 and 192\,px; and NanoDet-Plus-m~\cite{rangilyu2021nanodet} at 160\,px (0.218 GFLOPs), COCO-pretrained, of which we count only the inference parameters. \tbd{State the detector training recipe: classes, epochs, augmentation, how one hand box and one target box are selected per image, and the validation rule used to tune the hand-gate thresholds of the end-to-end score.} Every model is scored by the same evaluation code on the same images, and GFLOPs are counted with the same \texttt{torch} FLOP counter. Model sizes are the serialised SpatialNet state dicts and the exported ONNX files of the detectors.

\paragraph{Controlled detector runs.} To obtain paired intervals and to test two open questions, we also trained YOLO26n ourselves under one fully specified recipe (Ultralytics 8.4.177): the converted EgoPoint images at $128\times128$; two classes, the wearer's hand and the target, with no labels on no-hand images; 150 epochs with early stopping after 50 epochs without improvement, batch 64 and seed 42; and Ultralytics' default augmentation without mosaic, mixup or copy-paste. The backbone starts from ImageNet (the YOLO26n classification weights, which initialise its first ten layers, with the neck and head from scratch) or from COCO. One ImageNet-start run is also trained with our distractor pasting, applied on the fly with the same code and probabilities as for SpatialNet (Section~\ref{sec:ownership}); the pasted hands stay unlabelled background, since a detector has no ownership-weighted heatmap loss. For each image we score the top-scoring hand box and target box, and for every model, SpatialNet included, the hand-gate threshold is the value that maximises no-hand versus hand accuracy on the validation split.

\paragraph{Downstream recognition test.} IoU and centre in target describe the box, not what a later step can do with it (Section~\ref{sec:problem}). To measure that directly, we crop each model's predicted target box from the test image, restore the photograph's original aspect ratio, pad the crop to a square, and let an off-the-shelf zero-shot recogniser choose among the 348 target category names of EgoPoint, as a ``what is that?'' assistant would. We use two recognisers trained with different objectives, CLIP ViT-B/16~\cite{radford2021clip} and SigLIP base~\cite{zhai2023siglip} (16-pixel patches, 224\,px input), without any training, and average four prompt templates (``a photo of a \{\}.'' and three variants). A crop counts as recognised if the top-ranked name is the labelled category (top-1; we also check top-5). Crops of the labelled box and of the whole frame give an upper and a lower reference. Crops are taken from the 320\,px copy of each image. The test covers the 2,603 test targets with a category name and is category-level: a box on a neighbouring object of the same category counts as recognised, and near-synonyms such as ``bike'' and ``bicycle'' count as errors, for every model alike.

\paragraph{Pretraining overlap.} The COCO-pretrained weights are trained on COCO train2017, which contains every COCO train2014 photograph. These detectors have therefore been pretrained on the source photographs of at least 41\% of our test images (Section~\ref{sec:data}), together with box annotations of the objects they contain. Section~\ref{sec:detectors} tests whether this inflates their scores, by comparing COCO and ImageNet starts on COCO-derived and other test images. ImageNet pretraining, which our backbone and the second YOLO26n start use, and training from scratch involve no COCO detection data, so YOLO26n with an ImageNet start is the like-for-like reference for SpatialNet.

\paragraph{Near-duplicates.} 141 test images (4.7\%) are byte-identical to a training or validation image, and 362 have an exact or near copy there. We therefore repeat the headline comparisons on the test split without the 141 duplicates. Unless stated otherwise, every conclusion holds on that duplicate-free set; the main model's P@0.5 moves from 58.9\% to 57.2\%, inside its confidence interval. In addition, 230 test images share their COCO source photograph with a training or validation image; this overlap within EgoPoint affects all trained models alike.

\paragraph{Model selection.} The reported checkpoint has the lowest validation loss. Training, early stopping and checkpoint selection use only the training and validation splits; the test split is evaluated once, after training. The ablation, resolution and seed runs follow the same rule and do not change the reported model.
